% SPDX-License-Identifier: Apache-2.0
\section{The Reset Nothing Declares}
\label{sec:x-reset}

A clock is not a port. A unit waits for an edge of it, and the
netlist gives every module a \code{clk} and joins each child to its
parent's, because a clock reaches everything and a design that
threaded it by hand would say nothing by doing so.

The reset is the same, and for the same reason. No unit below
declares one, and every module the netlist writes has an \code{rst}
beside its \code{clk}. While it is high, every register goes back to
what it held before the first edge, and the channels between units
empty: what a channel holds is two values and two valid bits, and
clearing the bits is what makes it empty, so the data behind them is
left where it is and never read.

That is what a reset is for. A register put back is a register whose
value nobody has to reason about, and a design whose every register
is back where it started is the design that came out of the
configuration, which is the state that is known to work. Before this,
a reset reached the core, the timer, the serial port and the
interrupt controller, and reached nothing on the bus: a reset caught
while a transaction was in flight left a bridge waiting for an
acknowledgement that the reset had already cancelled, and the next
access to that memory waited behind it for ever (issue 346).

\code{Ticker} counts the ticks it is given, and says nothing about
any of this. The run counts four, holds the reset for two cycles with
the tick still high, and lets go.

\lstinputlisting[caption={\code{ex\_reset.rs}. Run with \code{bazel run //lib/examples:ex\_reset}.},label={lst:x-reset}]{lib/examples/ex_reset.rs}

\lstinputlisting[style=ascii,caption={What \code{ex\_reset} printed when the build ran it: the count, then zero while the reset is held although the tick is still high, then counting again from there.},label={lst:x-reset-out}]{out_reset.txt}

A unit that wants to \emph{do} something under reset, rather than
merely forget, takes \code{rst} as an \code{In<Bit>} and reads it.
The netlist joins that port to the same net instead of adding a
second of the same name, and adds no clearing of its own: a unit that
answered the reset itself has said what it means by it, and the Rust
and the netlist run that same answer. The core does this, because
forgetting where it was is not enough; it has to fetch from the
reset vector. \code{rst} is therefore a name the netlist keeps, and a
port of that name which is not an input is refused when the unit is
lowered, which is why the watchdog's request to be reset is called
\code{rst\_req}.

The reset is in the trace like any other signal, so the testbenches
drive it at the cycles the Rust run asserted it, and the VHDL under
\code{nvc} and the Verilog under Verilator are checked against the
Rust through the reset and out the other side.

\begin{figure}[htbp]
\centering
\input{reset_timing.tex}
\caption{\code{count} through a reset: four ticks counted,
\code{rst} high for two cycles with \code{tick} still high, and the
count away from zero again after it.}
\label{fig:x-reset}
\end{figure}
